\section{Verified candidates, physical margins and robustness}
\label{sec:orf-verified-followup}
This separately registered follow-up uses the same 24 within-family scenarios,
not new independent instances. Seeds 211, 257 and 293 replace the earlier seeds.
The 630 assigned MIP runs comprise 288 candidate-target runs, 126 additional
precision runs and 216 robustness runs. Reused controls are identified below.
One additional technical repeat standardizes the first Base run's verification
timing after an implementation amendment; its 54 candidate vectors, solver work,
endpoints and first-target indices reproduce exactly. Both traces, the amendment
and a worker failure before optimizer creation are retained. The repeat adds no
independent observation. The frozen research note describes the pre-integration
editorial assessment; the present manuscript incorporates only selected findings.

\subsection{Targets, verification and timing}
Gurobi 13.0.2 uses four threads, a 30-second limit, a requested 1\% MIP gap,
automatic presolve, \texttt{FeasibilityTol=1e-6},
\texttt{IntFeasTol=1e-9}, \texttt{NumericFocus=0} and automatic scaling, except
for the explicitly named controls. L2-NF3 changes only \texttt{NumericFocus} to 3.
Runs use fresh sequential processes on the declared local machine.
The reference physical requirements are 100\,m$^3$ for activation and water
balance, and 0.01\,MW for power.
All callback candidate events and the final postsolve point are saved;
callback events need not improve the incumbent. There are 6,725 primary
candidate events, with 1,943 full rational candidate audits. Exact rounded costs
reject the remaining candidates when neither fixed quality threshold is possible;
their full vectors remain available for replay.

For each scenario, $U_{\rm ref}$ is the exact cost of a rational feasible Base
witness acquired before these runs. A target uses
$C(\bar x)\le U_{\rm ref}+\tau\max\{1,|U_{\rm ref}|\}$, with primary
$\tau=0.001$ and secondary $\tau=0.01$. Neither reference is assumed optimal.
Qualification requires physical budgets at both raw and integer-rounded points,
raw integer distances at most $10^{-9}$, and exact variable bounds at the rounded
point. For rows without a declared physical budget, the rounded-point backward
error is at most $10^{-6}$, with denominator
$\max\{1,|b|,\sum_j|a_j\bar x_j|\}$.
This is a budget-verified approximate candidate, not exact feasibility of every
row or a MILP optimality certificate. The solver's requested 1\% stopping gap
can terminate before the primary 0.1\% target, without any numerical failure.

\emph{Discovery} is the elapsed fresh-process time to the first candidate later
found to qualify. \emph{Verified availability} is the actual elapsed time when
the postsolve rational audit reaches it, including solver completion, point
storage, parsing and earlier candidate checks. Neither is an online verified
stopping experiment. PAR2 retains every paired assignment: a hit contributes
its measured time, while a miss contributes 60 seconds plus measured process
preparation. A PAR2 improvement can therefore reflect more target hits rather
than faster common successes. Common-success ratios describe a selected subset
and are reported alongside all-assignment scores.

The uncertainty calculation uses 100,000 whole-scenario bootstrap samples,
retaining all three paired seeds within each sampled scenario. Marginal 95\%
percentile intervals are computed on log ratios; the primary analysis seed is
2026101906. The 24-, 12- and six-scenario analyses are conditional on known
parent structures, rather than uncertainty estimates for unrelated MIP families.
An independent implementation reconstructs all 44 reported intervals.
% Generated by scripts/research/orf_followup_tables.py
\begin{table}[htbp]
\centering
\small
\caption{Candidate-study final checks and target attainment; all 72 assignments per policy are retained.}
\label{tab:orf-verified-targets}
\resizebox{\linewidth}{!}{%
\begin{tabular}{lrrrrrr}
\toprule
Policy & Target 0.1\% & Target 1\% & Physical & Domains & Other rows & Bound flags \\
\midrule
Base & 53 & 68 & 72 & 69 & 72 & 2 \\
L2 & 51 & 66 & 72 & 72 & 72 & 7 \\
Budget-L2 & 54 & 72 & 72 & 72 & 72 & 1 \\
L2-NF3 & 54 & 72 & 72 & 72 & 72 & 0 \\
\bottomrule
\end{tabular}
}
\end{table}

% Generated by scripts/research/orf_followup_tables.py
\begin{table}[htbp]
\centering
\small
\caption{Budget-L2 relative to each control: paired geometric penalized scores and common-success times.}
\label{tab:orf-verified-par2}
\resizebox{\linewidth}{!}{%
\begin{tabular}{llrllr}
\toprule
Control & Target & Common pairs & Discovery PAR2 [95\%] & Availability PAR2 [95\%] & Common availability \\
\midrule
Base & 0.1\% & 45 & 0.772 [0.375, 1.596] & 0.802 [0.405, 1.595] & 0.742 \\
Base & 1\% & 68 & 0.613 [0.374, 0.893] & 0.638 [0.401, 0.901] & 0.816 \\
L2 & 0.1\% & 44 & 0.875 [0.422, 1.807] & 0.876 [0.441, 1.722] & 0.980 \\
L2 & 1\% & 66 & 0.759 [0.443, 1.124] & 0.753 [0.454, 1.085] & 1.041 \\
L2-NF3 & 0.1\% & 46 & 0.953 [0.497, 1.834] & 0.943 [0.509, 1.746] & 0.994 \\
L2-NF3 & 1\% & 72 & 1.014 [0.941, 1.093] & 1.010 [0.950, 1.082] & 1.010 \\
\bottomrule
\end{tabular}
}
\end{table}

All four policies satisfy physical budgets in 72/72 final points. Exact final
domain checks pass in 69/72 Base runs and all 72 runs of the other policies.
Remaining-row backward errors pass in every run. Budget-L2 still has one
independently refuted bound, versus seven for L2 and none for L2-NF3; primal
qualification does not remove this distinction. Bound contradictions use feasible
witnesses of the actual imported representation. An undetected contradiction is
not a certificate for every reported bound.

\subsection{Physical requirements and margins}
Six scenarios (indices 01 and 08 at each horizon) and three seeds test three
uniform water budgets of 10, 100 and 1,000\,m$^3$, with power budget 0.01\,MW.
A heterogeneous specification uses activation 10\,m$^3$, water-balance
0.001\,m$^3$ and power 0.001\,MW. These are numerical study requirements, not
operator standards; the one-liter balance is a numerical conservation goal,
not a measurement-accuracy assertion. Each row's physical-to-algebraic unit
conversion is applied before selection and reconstructed independently.
Budget-L2 at 100\,m$^3$ and L2 reuse candidate-study runs. The other three
Budget-L2 specifications add 54 runs; Export-Budget-L2 adds 72, for 126 new calls.
% Generated by scripts/research/orf_followup_tables.py
\begin{table}[htbp]
\centering
\small
\caption{Observed physical margins under four specifications; each policy passes all 18 final points per specification.}
\label{tab:orf-verified-precision}
\resizebox{\linewidth}{!}{%
\begin{tabular}{llrrr}
\toprule
Specification & Each policy passes & Max $q$, L2 & Max $q$, Budget-L2 & Max $q$, Export-Budget-L2 \\
\midrule
10\,m$^3$ & 18/18 & 6.44e-08 & 6.44e-08 & 7e-08 \\
100\,m$^3$ & 18/18 & 6.44e-09 & 6.77e-09 & 7e-09 \\
1,000\,m$^3$ & 18/18 & 6.44e-10 & 6.44e-10 & 6.34e-10 \\
Heterogeneous & 18/18 & 0.000644 & 0.000681 & 0.000712 \\
\bottomrule
\end{tabular}
}
\end{table}

Every policy passes all 18 final physical-budget checks in every specification.
No row has normalized residual $q$ near the limit (the recorded near-limit range
is $0.5\le q\le1$). Rechecking the reused 100\,m$^3$ Budget-L2 points under all
four requirements also passes. Thus this sample supplies margins and contract
interventions, without additional primal recovery over L2.

The independent uniform-reserve audit performs 95,040 row/configuration checks.
At the 100\,m$^3$ reference budget, 16,896 of 25,344 L2 physical rows satisfy the
full export-and-rounding sufficient condition; Budget-L2 satisfies it on all
25,344 rows across the 24 scenarios. The new selector certifies all 25,344
row--budget cases on six scenarios and four specifications. These are repeated
row/specification checks, not independent experimental units. In this sample,
the old Budget-L2 factors already satisfy the actual-export condition: the new
selector integrates the check into selection rather than recovering missing
conditional coverage. A compiled numerical-envelope control reproduces every
old factor exactly. Once the factors coincide, labels alone do not change the
solver's imported model.

A precheck at a 1\,m$^3$ activation budget finds 8,256 rows in 22/24 scenarios
with $E\ge\delta$ at integer distance $10^{-9}$. Because $\epsilon>0$ and
$B\ge0$, no factor can satisfy the stated sufficient contract for those rows.
This is neither dispatch infeasibility nor proof that an actual point cannot
meet 1\,m$^3$. No MIP is solved for this precheck, and the two other scenarios
are not thereby certified compatible in all respects.

\subsection{CPLEX and internal scaling}
CPLEX 22.1.2 tests 12 scenarios (indices 01, 03, 05 and 07 at each horizon),
three seeds and Base/L2/Budget-L2, adding 108 runs. This prespecified subset
contains three of the known parents. All policies use traditional MIP search
to expose incumbent callbacks, four threads, 30 seconds, requested gap 1\%,
primal tolerance $10^{-6}$, integer tolerance $10^{-9}$ and default scaling.
The full native library is identified by checksum; the default Community API's
size limit is documented by separate preflight probes. Imported matrices, RHS,
objectives, bounds and domains are checked against the intended exports.
Solver work units are compared only within the same solver.

Gurobi tests \texttt{ScaleFlag=0} and 2 on the six precision-study scenarios,
three seeds and the same three policies, adding 108 runs. Automatic-scaling
controls reuse 54 candidate-study runs. All final physical budgets pass, even
where the strict candidate criterion or independently checked dual bounds fail.
% Generated by scripts/research/orf_followup_tables.py
\begin{table}[htbp]
\centering
\small
\caption{Robustness subsets at the primary target; paired entries are Budget-L2 then L2. SF denotes ScaleFlag.}
\label{tab:orf-verified-robustness}
\resizebox{\linewidth}{!}{%
\begin{tabular}{lrlll}
\toprule
Configuration & Scenarios & Target hits, Budget; L2 & Availability PAR2 [95\%] & Bound flags \\
\midrule
CPLEX & 12 & 15/36; 15/36 & 1.021 [1.005, 1.040] & 5; 6 \\
Gurobi SF=-1 & 6 & 14/18; 14/18 & 0.934 [0.128, 4.635] & 0; 3 \\
Gurobi SF=0 & 6 & 15/18; 13/18 & 0.689 [0.098, 3.571] & 0; 4 \\
Gurobi SF=2 & 6 & 15/18; 14/18 & 0.782 [0.109, 4.010] & 0; 4 \\
\bottomrule
\end{tabular}
}
\par\smallskip
\begin{minipage}{\linewidth}\footnotesize
All policies satisfy final physical budgets. Exact-bound rejections affect CPLEX target counts; the post hoc projection sensitivity is reported in the text. Work units are not compared across solvers.
\end{minipage}
\end{table}

CPLEX's strict exact-bound rule rejects tiny final continuous-bound deviations
between approximately $1.3\times10^{-13}$ and $4.2\times10^{-11}$; these should
not be interpreted as substantial physical inaccuracies. A clearly post hoc
sensitivity rounds integers, projects onto exact variable bounds and rechecks
the transformed point's physical rows, backward errors and exact cost for every
policy. CPLEX then attains the primary target in 24/36 Base, 30/36 L2 and 30/36
Budget-L2 runs; all three attain the secondary target in 36/36. The largest
CPLEX candidate projection is approximately $2.2\times10^{-9}$.
The primary results and definition remain unchanged. No online availability
times are claimed for these projected points, and projection does not validate
the contradicted dual bounds. Thus CPLEX supplies a transfer/control result,
without an incremental Budget-L2 benefit over L2. The wide six-scenario
\texttt{ScaleFlag} intervals likewise establish no general advantage.

\subsection{Constructive selector and independent checks}
The export-aware policy uses \eqref{eq:orf-export-contract} during selection.
For every candidate it forms the products, writes 17 significant decimal
digits, parses binary64, and computes exact rational differences from $\alpha a$
and $\alpha b$. It rejects overflow, underflow of a nonzero coefficient, actual
written magnitudes outside $[10^{-9},10^9]$, a nonzero coefficient error without
a finite absolute variable bound, or a failed sufficient contract. Zero error
needs no variable bound. RHS products must also remain finite.

For $\delta>E$, the implementation first forms inward-rounded binary64
endpoints of the ideal search interval
$[\max\{\epsilon/(\delta-E),\ell/m\},U/M]$.
Its finite candidate list contains the constrained L2 target and endpoints,
32 adjacent representable neighbors in each direction from that target,
65 logarithmically spaced trial values, and every representable power of two
inside the interval. Duplicates are removed. Candidates are ordered by the
computed binary64 key
$(|\log\alpha-\log\alpha_{L2}|,-\alpha)$, so ties prefer the larger factor.
The ranked implementation returns the first exact certificate; an exhaustive
implementation checks the same list and returns its first successful key.
Consequently both implementations select the same factor under this numerical
ranking. This is not exact real-logarithm ordering or global optimization over
all binary64 factors. Restricting the ideal search interval may miss factors
whose written coefficients happen to satisfy the limits after rounding.

If $E\ge\delta$, the sufficient condition is empty for every positive factor.
Other empty-search or failed-list outcomes are explicit abstentions, not a
complete nonexistence proof for actual-export factors. For a returned factor,
the guarantee still requires the raw variable bounds, actual written residual
and integer distances used in its certificate. Nominal solver parameters do not
prove these hypotheses; direct original-point verification remains necessary.

Independent tests check 2,000 rational points on exported-tolerance boundaries
on 80 rows with nonzero coefficient export errors, compare ranked and exhaustive
selection, and exercise missing-bound and rounding-incompatibility cases.
A controlled example uses $a=(3)$, $b=0.1$, $X=10^{16}$ and $\delta=0.01$:
the ideal binary64 factor 0.3333333333333333 fails the uniform export reserve,
whereas the neighboring selected factor 0.33333333333333326 passes. This is a
software test, not evidence of an applied speedup or a second-family benefit.
The ranked implementation reproduces all 24 complete experimental LP exports
by SHA-256. A later selection-only benchmark measures medians of 0.630 seconds
for the ranked procedure and 8.175 seconds for the exhaustive prototype.
It compares implementations of the same selector, not Budget-L2 against L2;
it neither replaces recorded MIP process times nor establishes an end-to-end
acceleration. Inputs, all outcomes, failed attempts and reproducible checks are
supplied in Online Resource~2.
